\documentclass[UTF8,a4paper,11pt,oneside,fontset=none]{ctexrep}
\usepackage[margin=22mm,headheight=16pt]{geometry}
\setCJKmainfont{Noto Serif CJK SC}
\setCJKsansfont{Noto Sans CJK SC}
\setCJKmonofont{Noto Sans Mono CJK SC}
\setmainfont{TeX Gyre Pagella}
\setsansfont{TeX Gyre Heros}
\setmonofont{DejaVu Sans Mono}[Scale=0.86]
\usepackage{xcolor,longtable,booktabs,array,enumitem,fvextra,fancyhdr,xurl}
\usepackage[unicode,colorlinks=true,linkcolor=blue!45!black,urlcolor=blue!55!black,bookmarksnumbered=true]{hyperref}
\hypersetup{pdftitle={遥感算法容器封装与本机部署操作手册},pdfauthor={项目操作文档},pdfsubject={Docker、Kubernetes、GPU、离线交付与排错}}
\urlstyle{tt}
\DefineVerbatimEnvironment{Verbatim}{Verbatim}{fontsize=\small,breaklines=true,breakanywhere=true,frame=single,framesep=2.5mm,rulecolor=\color{black!20},commandchars=\\\{\}}
% Commands are literal; no TeX escape interpretation inside shell listings.
\RecustomVerbatimEnvironment{Verbatim}{Verbatim}{fontsize=\small,breaklines=true,breakanywhere=true,frame=single,framesep=2.5mm,rulecolor=\color{black!20}}
\setlist{itemsep=3pt,topsep=5pt,leftmargin=2em}
\setlength{\parskip}{5pt}
\setlength{\parindent}{2em}
\setlength{\tabcolsep}{3pt}
\renewcommand{\arraystretch}{1.22}
\emergencystretch=2em
\raggedbottom
\pagestyle{fancy}
\fancyhf{}
\fancyhead[L]{\small 遥感算法封装与部署}
\fancyhead[R]{\small 操作手册 · 2026-09-08}
\fancyfoot[C]{\thepage}
\fancypagestyle{plain}{\fancyhf{}\fancyfoot[C]{\thepage}\renewcommand{\headrulewidth}{0pt}}
\ctexset{chapter={format=\Huge\bfseries\sffamily,beforeskip=5pt,afterskip=22pt},section={format=\Large\bfseries\sffamily}}
\setcounter{tocdepth}{1}
\DefineVerbatimEnvironment{TerminalOutput}{Verbatim}{fontsize=\footnotesize,breaklines=true,breakanywhere=true,frame=leftline,framesep=2.5mm,rulecolor=\color{black!40}}
\DefineVerbatimEnvironment{Configuration}{Verbatim}{fontsize=\small,breaklines=true,breakanywhere=true,frame=single,framesep=2.5mm,rulecolor=\color{blue!40!black}}
\begin{document}
\begin{titlepage}
\vspace*{22mm}
{\sffamily\color{blue!45!black}\Large REMOTE SENSING · OPERATIONS\par}
\vspace{12mm}
{\sffamily\Huge\bfseries 遥感算法容器封装\par 与本机部署操作手册\par}
\vspace{10mm}
{\Large Docker / Kubernetes / GPU / 离线交付\par}
\vspace{14mm}
{\large 适用算法镜像：0.3.0-kmeans1\par 文档版次：1.1\par 编制日期：2026 年 9 月 8 日\par}
\vfill
\noindent\begin{minipage}{0.95\textwidth}
本手册规定遥感算法封装测试版的部署、配置、执行、停止、验证与交付操作。各操作条目包含输入填写要求、指令及参数、终端输出和完成判定，适用于当前本机测试环境及已准备基础依赖的 CPU 离线运行环境。\par
\vspace{5mm}
适用基线：算法镜像 0.3.0-kmeans1、源码提交 e59ce62。本文引用的实测输出来自 2026-09-08 已保存记录，动态字段不作为固定验收值。
\end{minipage}
\end{titlepage}
\pagenumbering{roman}
\tableofcontents
\clearpage
\pagenumbering{arabic}
\chapter{文档约定与操作流程}
\section{适用范围与版本}
本手册适用于 Linux amd64 上的遥感算法容器封装测试版。已接入算法为 MNDWI、KMEANS；smoke 用于文件流转诊断。GPU 章节验证设备接入与 CUDA 运算，不改变上述 CPU 示例的计算后端。当前 Kubernetes 实现为本机单节点 kind 集群 rs-lab。

环境和测试基线来自同日保存的部署记录；本次文档修订不安装环境、不运行运维命令、不合并业务脚本。正式目标平台、真实算法、网络与安全门禁须按实际交付要求验收。

\section{指令与输出表示规则}
\begin{center}\small
\begin{tabular}{>{\raggedright\arraybackslash}p{4.84cm}>{\raggedright\arraybackslash}p{10.76cm}}
\toprule
\textbf{表示方式} & \textbf{含义} \\ \midrule
执行命令 & 复制到 Bash 终端执行；不包含 \$ 提示符 \\
配置文件内容 & 保存到指定 env 文件；不是直接执行的 shell 命令 \\
需填写内容 & 使用前应替换或核对的变量与路径 \\
终端输出（实测记录） & 从已保存证据中摘录；保留来源路径 \\
终端输出（格式示意） & 按现有脚本行为构造的显示示例；不宣称重新实测 \\
无输出 & 命令成功时可能不打印文字；应查看退出码或后续状态 \\
\bottomrule
\end{tabular}
\end{center}
同一操作条目的变量在当前终端有效。跨终端操作必须重新填写容器名或 Job 名。代码块中的反斜杠是 Bash 续行符；复制命令推荐使用配套 Markdown 文本版，避免 PDF 的视觉折行进入命令。

\section{标准操作顺序}
\begin{center}\small
\begin{tabular}{>{\raggedright\arraybackslash}p{2.5cm}>{\raggedright\arraybackslash}p{7.18cm}>{\raggedright\arraybackslash}p{5.93cm}}
\toprule
\textbf{阶段} & \textbf{操作} & \textbf{形成的记录} \\ \midrule
准备 & 确认项目路径、数据盘、镜像及访问权限 & 环境检查输出 \\
配置 & 建立独立任务 env，填写输入、输出和算法参数 & env 配置文件 \\
执行 & 启动、保存 \path|stdout/stderr|、记录退出码 & 日志与退出码文件 \\
判定 & 检查 result 文件与 success.json & 结果及完成标记 \\
停止或恢复 & 停止指定任务；失败时采用新目录重跑 & 中止记录或新任务结果 \\
验收与交付 & 执行相应测试、扫描、导出和校验 & 验收报告、镜像及 SHA256 \\
\bottomrule
\end{tabular}
\end{center}
\chapter{环境初始化与状态确认}
\section{ENV-01 设置操作终端}
\par\noindent\textbf{用途}\par
设定项目目录、算法镜像和 Kubernetes 工具路径。

\par\noindent\textbf{执行前提}\par
本机目录已存在；使用有 Docker 权限的普通用户，日常运行不使用 sudo。

\par\noindent\textbf{需填写内容}\par
\path|PROJECT_ROOT|：项目根目录；IMAGE：准备运行的已加载镜像。本机示例可直接使用下列值。

\par\noindent\textbf{执行命令}\par
\begin{Verbatim}
export PROJECT_ROOT=/mnt/robot_disk/remote_sensing_project
export IMAGE=remote-sensing-framework:0.3.0-kmeans1
export PATH=/mnt/robot_disk/k8s/bin:$PATH
export KUBECONFIG=/mnt/robot_disk/k8s/config/kubeconfig
cd "$PROJECT_ROOT"
printf 'PROJECT_ROOT=%s\nIMAGE=%s\n' "$PWD" "$IMAGE"
\end{Verbatim}
\begin{center}\small
\begin{tabular}{>{\raggedright\arraybackslash}p{5.3cm}>{\raggedright\arraybackslash}p{10.3cm}}
\toprule
\textbf{指令 / 变量} & \textbf{说明} \\ \midrule
export & 将变量传递给后续子进程；重新开终端后需再次设置 \\
cd "\$\path|PROJECT_ROOT|" & 切换工作目录；引号支持路径中的空格 \\
PATH & 追加本机 kind 和 kubectl 所在目录 \\
KUBECONFIG & 指定集群访问配置；不得作为公开交付附件 \\
IMAGE & 选择算法镜像；不触发拉取、构建或 GPU 启用 \\
\bottomrule
\end{tabular}
\end{center}
\par\noindent\textbf{终端输出（格式示意）}\par
\begin{TerminalOutput}
PROJECT_ROOT=/mnt/robot_disk/remote_sensing_project
IMAGE=remote-sensing-framework:0.3.0-kmeans1
\end{TerminalOutput}
示意中的任务名、时间、目录、哈希和耗时以实际运行值为准；该输出用于说明格式，不构成本次重新运行证据。

\par\noindent\textbf{完成判定与异常处理}\par
打印的项目路径和镜像名与计划一致。cd 失败时先修正 \path|PROJECT_ROOT|；不要在其他工作目录继续执行相对路径命令。

\section{ENV-02 检查磁盘和 Docker}
\par\noindent\textbf{用途}\par
确认数据盘挂载、Docker 服务、存储目录和镜像可用。

\par\noindent\textbf{执行前提}\par
已完成 ENV-01；本机 GPU 和 Kubernetes 检查不作为 CPU 运行的必需前提。

\par\noindent\textbf{需填写内容}\par
无额外必填项；目标机器应将数据盘路径替换为已确认的实际路径。

\par\noindent\textbf{执行命令}\par
\begin{Verbatim}
findmnt -T "$PROJECT_ROOT"
df -h /mnt/robot_disk
systemctl is-active docker
docker info --format 'DockerRootDir={{.DockerRootDir}}'
docker image inspect "$IMAGE" --format '{{.Id}}'
bash scripts/check-environment.sh .
\end{Verbatim}
\begin{center}\small
\begin{tabular}{>{\raggedright\arraybackslash}p{5.46cm}>{\raggedright\arraybackslash}p{10.14cm}}
\toprule
\textbf{参数} & \textbf{说明} \\ \midrule
findmnt -T 路径 & 显示承载该路径的文件系统 \\
df -h & 按便于阅读的单位显示容量和剩余空间 \\
--format & 提取指定字段，避免输出整份配置 \\
check-environment.sh . & 检查工具与 CPU 基础条件；句点用于显示当前目录所在盘 \\
\bottomrule
\end{tabular}
\end{center}
\par\noindent\textbf{终端输出（格式示意）}\par
\begin{TerminalOutput}
active
DockerRootDir=/mnt/robot_disk/docker-engine
sha256:e9bcee28fc600022a4a919897528e703e1c30d14730b32679c57a4c9a7ea3c89
CPU_PREREQUISITE_EXIT=0
\end{TerminalOutput}
示意中的任务名、时间、目录、哈希和耗时以实际运行值为准；该输出用于说明格式，不构成本次重新运行证据。

\par\noindent\textbf{完成判定与异常处理}\par
Docker 为 active，CPU 基础检查为 0，镜像 inspect 成功。本机 DockerRootDir 应落在 /\path|mnt/robot_disk| 下。\path|GPU_CHECK=NOT_READY| 不代表 CPU 不可运行，但应停止进入 GPU 验收流程。

\section{当前存储映射}
\begin{center}\small
\begin{tabular}{>{\raggedright\arraybackslash}p{5.62cm}>{\raggedright\arraybackslash}p{4.52cm}>{\raggedright\arraybackslash}p{5.46cm}}
\toprule
\textbf{用户填写的宿主目录} & \textbf{容器内固定目录} & \textbf{读写用途} \\ \midrule
\path|INPUT_DIR| & /\path|data/input| & 只读输入 \\
\path|OUTPUT_DIR| & /\path|data/output| & 持久输出；新任务使用新目录 \\
\path|WORK_DIR| & /\path|data/work| & 大型临时文件预留，现有两个示例未使用 \\
项目 \path|data/input| & kind 内 /rs-input & Pod 输入来源 \\
项目 outputs & kind 内 /rs-output & Pod 结果来源 \\
\bottomrule
\end{tabular}
\end{center}
本机数据盘为 /\path|dev/sdb2|，ext4 挂载于 /\path|mnt/robot_disk|；记录时可用约 858 GiB。Docker 数据目录为 /\path|mnt/robot_disk/docker-engine|，宿主 containerd 数据目录为 /\path|mnt/robot_disk/containerd|。程序和部分系统服务文件仍位于系统分区。

\chapter{任务配置文件规范}
\section{CFG-01 建立本次任务 env}
\par\noindent\textbf{用途}\par
将运行配置保存为独立文件，避免不同任务相互覆盖。

\par\noindent\textbf{执行前提}\par
执行 ENV-01；任务输出目录尚未包含已完成结果。

\par\noindent\textbf{需填写内容}\par
\path|TASK_ID|：本次任务标识。示例自动生成；正式使用可采用工单号加时间。生成文件后须核对 \path|INPUT_DIR|、\path|RS_INPUT| 和 \path|RS_PARAMS|。

\par\noindent\textbf{执行命令}\par
\begin{Verbatim}
TASK_ID="kmeans-$(date +%Y%m%d-%H%M%S)-$$"
ENV_FILE="$PROJECT_ROOT/.env.$TASK_ID"
cat > "$ENV_FILE" <<EOF
IMAGE=$IMAGE
INPUT_DIR=./data/input/kmeans-example
OUTPUT_DIR=./outputs/$TASK_ID
WORK_DIR=./work/$TASK_ID
CONTAINER_NAME=rs-$TASK_ID
CPU_LIMIT=1
MEMORY_LIMIT=1g
STOP_TIMEOUT=10
RS_ALGORITHM=KMEANS
RS_INPUT=/data/input/kmeans-input.tif
RS_PARAMS={"clusters":3,"seed":0,"max_iter":100,"max_pixels":1000000}
RS_TASK_ID=$TASK_ID
EOF
printf 'ENV_FILE=%s\nTASK_ID=%s\n' "$ENV_FILE" "$TASK_ID"
\end{Verbatim}
\par\noindent\textbf{终端输出（格式示意）}\par
\begin{TerminalOutput}
ENV_FILE=/mnt/robot_disk/remote_sensing_project/.env.kmeans-20260908-140000-12345
TASK_ID=kmeans-20260908-140000-12345
\end{TerminalOutput}
示意中的任务名、时间、目录、哈希和耗时以实际运行值为准；该输出用于说明格式，不构成本次重新运行证据。

\par\noindent\textbf{完成判定与异常处理}\par
env 文件存在，\path|INPUT_DIR| 指向真实输入目录，\path|RS_INPUT| 是容器内输入文件路径。模板中的 EOF 用于在创建时展开受控变量，env-run.py 实际读取文件时不执行 shell。

\section{需由使用者填写的字段}
\begin{center}\small
\begin{tabular}{>{\raggedright\arraybackslash}p{4.06cm}>{\raggedright\arraybackslash}p{6.08cm}>{\raggedright\arraybackslash}p{5.46cm}}
\toprule
\textbf{字段} & \textbf{填写规则} & \textbf{示例} \\ \midrule
\path|INPUT_DIR| & 必填；宿主目录 & ./\path|data/input/kmeans-example| \\
\path|OUTPUT_DIR| & 必填；每次新的宿主输出目录 & ./\path|outputs/kmeans-001| \\
\path|RS_ALGORITHM| & 选定已接入算法，区分大小写 & KMEANS \\
\path|RS_INPUT| & 必填；容器内文件名 & /\path|data/input/kmeans-input.tif| \\
\path|RS_PARAMS| & JSON 对象；内容见算法参数表 & \{"clusters":3,"seed":0\} \\
\path|CONTAINER_NAME| & 建议填写；当前任务唯一容器名 & rs-kmeans-001 \\
\path|RS_TASK_ID| & 建议填写；日志关联标识 & kmeans-001 \\
\bottomrule
\end{tabular}
\end{center}
\section{环境与资源字段}
\begin{center}\small
\begin{tabular}{>{\raggedright\arraybackslash}p{4.21cm}>{\raggedright\arraybackslash}p{6.4cm}>{\raggedright\arraybackslash}p{4.99cm}}
\toprule
\textbf{字段} & \textbf{模板值 / 缺省行为} & \textbf{设置影响} \\ \midrule
IMAGE & 模板为 0.3.0-kmeans1；运行脚本要求提供 & 选择已加载镜像 \\
\path|WORK_DIR| & 模板独立目录；未设置则 \path|OUTPUT_DIR/work| & 持久工作空间 \\
\path|CPU_LIMIT| & 1 & CPU 配额上限 \\
\path|MEMORY_LIMIT| & 模板 1g；run.sh 缺省 512m & 普通内存限制，不是显存 \\
\path|STOP_TIMEOUT| & 10；非负整数 & 正常终止宽限秒数 \\
\bottomrule
\end{tabular}
\end{center}
\path|INPUT_DIR|、\path|OUTPUT_DIR|、\path|WORK_DIR| 的相对路径以 env 文件所在目录为基准。文件支持中文和空格路径，不支持变量展开、命令替换及行尾注释。注释须独占一行；重复或未知字段被拒绝。env 覆盖继承的同名环境变量；追加的 CLI 算法参数覆盖对应运行参数。

\section{配置文件的填写样式}
\par\noindent\textbf{配置文件内容：自有数据示例（填写后保存，不在终端直接执行）}\par
\begin{Configuration}
IMAGE=remote-sensing-framework:0.3.0-kmeans1
INPUT_DIR=/mnt/robot_disk/robot_data/待处理影像
OUTPUT_DIR=/mnt/robot_disk/robot_data/结果/task-001
WORK_DIR=/mnt/robot_disk/robot_data/工作/task-001
CONTAINER_NAME=rs-task-001
CPU_LIMIT=1
MEMORY_LIMIT=1g
STOP_TIMEOUT=10
RS_ALGORITHM=KMEANS
RS_INPUT=/data/input/scene.tif
RS_PARAMS={"clusters":3,"seed":0,"max_iter":100,"max_pixels":1000000}
RS_TASK_ID=task-001
\end{Configuration}
需将示例目录和 scene.tif 替换为实际数据。当前输入波段使用原始数值，不自动标准化；正式数据的波段含义、缩放及聚类意义应另行确认。默认 env 不需要服务 URL、密码或访问令牌。

\chapter{算法执行与参数参考}
\section{RUN-01 生成 KMEANS 验收样本}
\par\noindent\textbf{用途}\par
生成确定性的 90×90 双波段合成 GeoTIFF，用于封装验收。

\par\noindent\textbf{执行前提}\par
已加载 IMAGE；仅在专用测试目录写入。已有同名样本可跳过生成。

\par\noindent\textbf{需填写内容}\par
样本目录：示例固定 \path|data/input/kmeans-example|。正式影像无需执行此步骤。

\par\noindent\textbf{执行命令}\par
\begin{Verbatim}
mkdir -p data/input/kmeans-example
docker run --rm --network none \
  --user "$(id -u):$(id -g)" \
  -v "$PWD/data/input/kmeans-example:/fixture" \
  -v "$PWD/tests:/tests:ro" \
  --entrypoint python "$IMAGE" \
  /tests/create_kmeans_fixture.py /fixture
rc=$?
printf 'fixture_exit=%s\n' "$rc"
ls data/input/kmeans-example
\end{Verbatim}
\begin{center}\small
\begin{tabular}{>{\raggedright\arraybackslash}p{5.93cm}>{\raggedright\arraybackslash}p{9.67cm}}
\toprule
\textbf{参数} & \textbf{用途} \\ \midrule
--rm & 退出后删除临时容器 \\
--network none & 此容器禁网，生成过程使用已有依赖 \\
--user UID:GID & 以宿主普通用户身份创建文件 \\
-v 宿主:容器[:ro] & 挂载样本目录及只读测试脚本 \\
--entrypoint python & 运行样本生成脚本，替换算法入口 \\
/fixture [side] & 输出目录及可选边长；缺省边长 90 \\
\bottomrule
\end{tabular}
\end{center}
\par\noindent\textbf{终端输出（格式示意）}\par
\begin{TerminalOutput}
fixture_exit=0
expected-labels.npy
kmeans-input.tif
\end{TerminalOutput}
示意中的任务名、时间、目录、哈希和耗时以实际运行值为准；该输出用于说明格式，不构成本次重新运行证据。

\par\noindent\textbf{完成判定与异常处理}\par
退出码 0，两个样本文件存在。生成脚本自身通常无 stdout；这里的状态文字来自 printf 和 ls。重新生成会覆盖指定测试样本，不得指向生产输入目录。

\section{RUN-02 执行 env 配置任务}
\par\noindent\textbf{用途}\par
运行指定算法并保存 stdout、JSON 日志及退出码。

\par\noindent\textbf{执行前提}\par
已完成 CFG-01 和输入准备；本终端保留 \path|ENV_FILE| 与 \path|TASK_ID|。

\par\noindent\textbf{需填写内容}\par
\path|ENV_FILE|：实际 env 文件；\path|TASK_ID|：用于宿主日志命名。若重新开终端，须手动恢复两者。

\par\noindent\textbf{执行命令}\par
\begin{Verbatim}
mkdir -p logs
if python3 scripts/env-run.py "$ENV_FILE" \
  > "logs/$TASK_ID.stdout" 2> "logs/$TASK_ID.jsonl"; then
  rc=0
else
  rc=$?
fi
printf '%s\n' "$rc" > "logs/$TASK_ID.exit"
printf 'task=%s exit=%s\n' "$TASK_ID" "$rc"
tail -n 2 "logs/$TASK_ID.jsonl"
\end{Verbatim}
\begin{center}\small
\begin{tabular}{>{\raggedright\arraybackslash}p{5.46cm}>{\raggedright\arraybackslash}p{10.14cm}}
\toprule
\textbf{指令 / 参数} & \textbf{说明} \\ \midrule
env-run.py \path|ENV_FILE| & 按字面解析 env，传给 run.sh；不执行配置文本 \\
> stdout 文件 & 保存标准输出；正常批处理时可能为空 \\
2> jsonl 文件 & 保存标准错误通道上的逐行 JSON 日志 \\
if ...; then ... else & 无论成功失败都记录真实退出码 \\
tail -n 2 & 显示日志最后两行；长任务可能有更多日志 \\
\bottomrule
\end{tabular}
\end{center}
\par\noindent\textbf{终端输出（格式示意）}\par
\begin{TerminalOutput}
task=kmeans-20260908-140000-12345 exit=0
{"level":"INFO","message":"algorithm_started","algorithm":"KMEANS", ...}
{"level":"INFO","message":"algorithm_completed","device":"cpu", ...}
\end{TerminalOutput}
示意中的任务名、时间、目录、哈希和耗时以实际运行值为准；该输出用于说明格式，不构成本次重新运行证据。

上面的省略号仅表示日志展示节选，不是完整合法 JSON。完整日志保留在 jsonl 文件中；本机已保存的完成日志含 \path|valid_pixels=8099|、\path|invalid_pixels=1|、iterations=2，来源 \path|reports/deployment/docker-integration/success.jsonl|。

\par\noindent\textbf{完成判定与异常处理}\par
退出码为 0 且出现 \path|algorithm_completed|。仍须执行 RUN-03 检查完成标记；单独看到启动日志不表示成功。退出非零时保留日志与输出目录，按故障章节定位。

\section{RUN-03 检查结果和完成标记}
\par\noindent\textbf{用途}\par
校验预期结果目录与任务状态。

\par\noindent\textbf{执行前提}\par
RUN-02 已结束。

\par\noindent\textbf{需填写内容}\par
\path|RESULT_DIR|：填写 env 中 \path|OUTPUT_DIR| 下的 result 目录；下列值适用于 CFG-01 生成的配置。

\par\noindent\textbf{执行命令}\par
\begin{Verbatim}
RESULT_DIR="$PROJECT_ROOT/outputs/$TASK_ID/result"
ls "$RESULT_DIR"
python3 - "$RESULT_DIR/success.json" <<'PYTHON'
import json, sys
from pathlib import Path
r = json.loads(Path(sys.argv[1]).read_text())
print("status=" + r["status"])
print("algorithm=" + r["algorithm"])
print("output=" + r["output"])
print("valid_pixels=" + str(r["valid_pixels"]))
assert r["status"] == "completed"
PYTHON
\end{Verbatim}
\par\noindent\textbf{终端输出（格式示意）}\par
\begin{TerminalOutput}
clusters.tif
success.json
status=completed
algorithm=KMEANS
output=/data/output/result/clusters.tif
valid_pixels=8099
\end{TerminalOutput}
示意中的任务名、时间、目录、哈希和耗时以实际运行值为准；该输出用于说明格式，不构成本次重新运行证据。

\par\noindent\textbf{完成判定与异常处理}\par
任务退出码 0、状态 completed、期望 GeoTIFF 存在共同构成运行完成判定。output 字段是容器内路径；宿主对应 \path|RESULT_DIR|。8099 仅适用于该合成样本，不是任意影像的固定结果。

\section{KMEANS 参数表}
\begin{center}\small
\begin{tabular}{>{\raggedright\arraybackslash}p{4.21cm}>{\raggedright\arraybackslash}p{3.59cm}>{\raggedright\arraybackslash}p{7.8cm}}
\toprule
\textbf{参数} & \textbf{类型 / 默认} & \textbf{允许值及含义} \\ \midrule
clusters & 整数 / 3 & 2 至 255；聚类数；有效且不同像元数不得小于该值 \\
seed & 整数 / 0 & 0 至 4294967295；随机种子 \\
\path|max_iter| & 整数 / 100 & 1 至 10000；每次初始化的最大迭代次数 \\
\path|max_pixels| & 整数 / 1000000 & 1 至 10000000；允许的影像总像元数上限 \\
\bottomrule
\end{tabular}
\end{center}
布尔值不作为整数接收；未知参数报错。输入最多 32 波段，全部波段参与计算，NoData 和非有限像元被排除。固定 \path|n_init=10|、单线程、Lloyd 算法；输出 uint16，0 为 NoData，1 至 K 为任意簇编号。允许的像元上限不保证对应内存足够；跨实现比较应考虑簇编号置换。

\section{RUN-04 执行 MNDWI 示例}
\par\noindent\textbf{用途}\par
计算绿光与 SWIR1 波段的 MNDWI，演示 CLI 参数覆盖。

\par\noindent\textbf{执行前提}\par
执行 ENV-01，存在镜像；目标输入为至少两个波段的专用测试 GeoTIFF。

\par\noindent\textbf{需填写内容}\par
本示例自动生成测试输入。自有影像应修改 \path|INPUT_DIR|、--input 以及两个波段编号。

\par\noindent\textbf{执行命令}\par
\begin{Verbatim}
mkdir -p data/input/mndwi-example logs
docker run --rm --network none \
  --user "$(id -u):$(id -g)" \
  -v "$PWD/data/input/mndwi-example:/fixture" \
  -v "$PWD/tests:/tests:ro" \
  --entrypoint python "$IMAGE" /tests/create_mndwi_fixture.py
TASK_ID="mndwi-$(date +%Y%m%d-%H%M%S)-$$"
if INPUT_DIR="$PWD/data/input/mndwi-example" \
  OUTPUT_DIR="$PWD/outputs/$TASK_ID" \
  WORK_DIR="$PWD/work/$TASK_ID" \
  CONTAINER_NAME="rs-$TASK_ID" RS_TASK_ID="$TASK_ID" \
  MEMORY_LIMIT=1g bash scripts/run.sh \
  --algorithm MNDWI --input /data/input/mndwi-input.tif \
  --params '{"green_band":1,"swir1_band":2,"output_name":"mndwi.tif"}' \
  > "logs/$TASK_ID.stdout" 2> "logs/$TASK_ID.jsonl"; then
  rc=0
else
  rc=$?
fi
printf '%s\n' "$rc" > "logs/$TASK_ID.exit"
printf 'task=%s exit=%s\n' "$TASK_ID" "$rc"
ls "outputs/$TASK_ID/result"
\end{Verbatim}
\par\noindent\textbf{终端输出（格式示意）}\par
\begin{TerminalOutput}
task=mndwi-20260908-141000-12345 exit=0
mndwi.tif
success.json
\end{TerminalOutput}
示意中的任务名、时间、目录、哈希和耗时以实际运行值为准；该输出用于说明格式，不构成本次重新运行证据。

\par\noindent\textbf{完成判定与异常处理}\par
退出码 0，结果文件与 success.json 存在。样本的数值校验由现有算法测试覆盖；真实影像另需确认反射率缩放与波段语义。

\section{MNDWI 参数表}
\begin{center}\small
\begin{tabular}{>{\raggedright\arraybackslash}p{4.21cm}>{\raggedright\arraybackslash}p{4.21cm}>{\raggedright\arraybackslash}p{7.18cm}}
\toprule
\textbf{参数} & \textbf{类型 / 默认} & \textbf{填写规则} \\ \midrule
\path|green_band| & 整数 / 1 & 从 1 开始，不能超过输入波段总数 \\
\path|swir1_band| & 整数 / 2 & 从 1 开始，不能超过输入波段总数 \\
\path|output_name| & 字符串 / mndwi.tif & 普通文件名，后缀 .tif 或 .tiff；不接受目录路径 \\
\bottomrule
\end{tabular}
\end{center}
计算公式为 (G-S1)/(G+S1)。输出 float32，分母为零或输入无效的像元为 NaN；保留空间参考。当前按窗口处理，不提供网络服务。

\section{算法命令行接口}
\begin{center}\small
\begin{tabular}{>{\raggedright\arraybackslash}p{4.68cm}>{\raggedright\arraybackslash}p{7.02cm}>{\raggedright\arraybackslash}p{3.9cm}}
\toprule
\textbf{参数} & \textbf{用途} & \textbf{对应配置} \\ \midrule
--algorithm & 选择算法 & \path|RS_ALGORITHM| \\
--input & 容器内输入文件 & \path|RS_INPUT| \\
--output-dir & 容器内输出目录；标准运行保持 /\path|data/output| & 由 run.sh 挂载决定 \\
--params & 算法 JSON 对象，外层用单引号包裹 & \path|RS_PARAMS| \\
--task-id & 日志关联标识 & \path|RS_TASK_ID| \\
--help / --list / --version & 帮助、接入状态、版本；不执行算法 & 无 \\
--healthcheck & 只检查 Python 和入口 & 不证明算法或挂载健康 \\
\bottomrule
\end{tabular}
\end{center}
\par\noindent\textbf{执行命令：查看入口说明}\par
\begin{Verbatim}
docker run --rm --network none "$IMAGE" --help
docker run --rm --network none "$IMAGE" --version
\end{Verbatim}
\par\noindent\textbf{终端输出（格式示意）}\par
\begin{TerminalOutput}
usage: ... [--algorithm ALGORITHM] [--input INPUT] ...
0.3.0
\end{TerminalOutput}
示意中的任务名、时间、目录、哈希和耗时以实际运行值为准；该输出用于说明格式，不构成本次重新运行证据。

帮助输出中的选项按镜像版本显示；--list 中 \path|not_implemented| 的算法不可作为已交付能力使用。

\chapter{运行监控、停止和重跑}
\section{CTL-01 查询状态与日志}
\par\noindent\textbf{用途}\par
定位运行中的容器并查看进展。

\par\noindent\textbf{执行前提}\par
任务仍在运行；run.sh 使用 --rm，小任务结束后容器会自动删除。

\par\noindent\textbf{需填写内容}\par
CONTAINER：填写 env 的 \path|CONTAINER_NAME|；示例值必须改为实际运行名称。

\par\noindent\textbf{执行命令}\par
\begin{Verbatim}
CONTAINER=rs-kmeans-001
docker ps --format 'table {{.Names}}\t{{.Status}}\t{{.Image}}'
docker logs --tail 20 "$CONTAINER"
docker stats --no-stream "$CONTAINER"
\end{Verbatim}
\begin{center}\small
\begin{tabular}{>{\raggedright\arraybackslash}p{4.99cm}>{\raggedright\arraybackslash}p{10.61cm}}
\toprule
\textbf{参数} & \textbf{用途} \\ \midrule
--tail 20 & 只显示最近 20 条容器日志 \\
--no-stream & 获取一次资源快照后退出 \\
docker logs -f & 可选持续跟踪；Ctrl+C 只结束查看 \\
CONTAINER & Docker 任务名，不是 \path|RS_TASK_ID| 或镜像标签 \\
\bottomrule
\end{tabular}
\end{center}
\par\noindent\textbf{终端输出（格式示意）}\par
\begin{TerminalOutput}
NAMES           STATUS          IMAGE
rs-kmeans-001   Up 12 seconds   remote-sensing-framework:0.3.0-kmeans1
CONTAINER ID   NAME            CPU %   MEM USAGE / LIMIT
<id>           rs-kmeans-001   ...     ...
\end{TerminalOutput}
示意中的任务名、时间、目录、哈希和耗时以实际运行值为准；该输出用于说明格式，不构成本次重新运行证据。

\par\noindent\textbf{完成判定与异常处理}\par
列表中的名字与目标一致。找不到容器时先查看宿主 logs 下的退出码；No such container 可能是 --rm 正常删除，而不是未运行。CPU、内存快照不是性能 SLA。

\section{CTL-02 停止指定容器}
\par\noindent\textbf{用途}\par
向指定任务请求正常终止，超时后由 Docker 强制终止。

\par\noindent\textbf{执行前提}\par
已确认容器名称，且没有将日志查看进程误认为算法进程。

\par\noindent\textbf{需填写内容}\par
CONTAINER：目标容器；\path|GRACE_SECONDS|：非负整数，默认 10 秒。

\par\noindent\textbf{执行命令}\par
\begin{Verbatim}
CONTAINER=rs-kmeans-001
GRACE_SECONDS=10
bash scripts/stop.sh "$CONTAINER" "$GRACE_SECONDS"
printf 'stop_command_exit=%s\n' "$?"
\end{Verbatim}
\par\noindent\textbf{终端输出（格式示意）}\par
\begin{TerminalOutput}
rs-kmeans-001
stop_command_exit=0
\end{TerminalOutput}
示意中的任务名、时间、目录、哈希和耗时以实际运行值为准；该输出用于说明格式，不构成本次重新运行证据。

\par\noindent\textbf{完成判定与异常处理}\par
\path|stop_command_exit=0| 仅表示 Docker 停止命令成功。算法运行终端记录的退出码可能为 143，强杀时常见 137；两者不得混淆。原生计算库可能延迟处理信号，当前无断点续跑保证。

\section{重跑规则}
\begin{itemize}
\item 正常协作取消会尝试清理临时目录；OOM、强杀或断电可能留下 .partial-* 与 .rs-lock。
\item 重跑时建立新的任务标识、容器名和 \path|OUTPUT_DIR|。保留旧配置、日志和结果目录，不在活动任务中清锁。
\item 只有运行退出码与完整 success.json 均满足要求才标记完成。停止与发布结果发生竞争时记录两者，不将非零退出直接作为成功。
\end{itemize}
\section{退出码参考}
\begin{center}\small
\begin{tabular}{>{\raggedright\arraybackslash}p{3.43cm}>{\raggedright\arraybackslash}p{12.17cm}}
\toprule
\textbf{退出码} & \textbf{含义} \\ \midrule
0 & 操作成功；信息命令不代表算法执行成功 \\
1 & 未处理错误 \\
2 & 参数、配置、输入内容或规模限制错误 \\
3 & 算法未知或未接入 \\
4 & 输入不存在/不可读，或算法输出 \path|I/O| 错误 \\
5 & 结果已存在、目录占用，或 smoke 输出失败 \\
130 / 143 & SIGINT / SIGTERM 协作中止 \\
137 & 强杀或可能 OOM；需结合平台状态 \\
125 / 126 / 127 & Docker 启动或入口执行错误 \\
\bottomrule
\end{tabular}
\end{center}
\chapter{CPU 与 Docker GPU 验收}
\section{TEST-01 CPU 封装回归}
\par\noindent\textbf{用途}\par
分别检查接口、容器运行约束和完整部署行为。

\par\noindent\textbf{执行前提}\par
设置 IMAGE；普通用户执行；完整验收会占用资源并创建诊断容器和输出。

\par\noindent\textbf{需填写内容}\par
IMAGE：待验收镜像；现有测试入口无额外必填 CLI 参数。

\par\noindent\textbf{执行命令}\par
\begin{Verbatim}
bash scripts/test-local.sh
bash scripts/test-container.sh
python3 tests/verify_deployment.py
\end{Verbatim}
\begin{center}\small
\begin{tabular}{>{\raggedright\arraybackslash}p{4.68cm}>{\raggedright\arraybackslash}p{5.93cm}>{\raggedright\arraybackslash}p{4.99cm}}
\toprule
\textbf{命令} & \textbf{检查范围} & \textbf{记录} \\ \midrule
test-local.sh & 接口与 env；算法库缺失可能跳过 & 终端 unittest 输出 \\
test-container.sh & smoke、错误输入、用户和挂载权限 & \path|outputs/container-test.*| \\
\path|verify_deployment.py| & 算法、中文路径、env、权限、OOM、停止重跑 & \path|outputs/deployment-*| \\
\bottomrule
\end{tabular}
\end{center}
\par\noindent\textbf{终端输出（实测记录）}\par
\begin{TerminalOutput}
default UID verified
filesystem restrictions verified
PASS: 框架容器检查通过；不是遥感算法验收。记录目录：/mnt/robot_disk/remote_sensing_project/outputs/container-test.cyuHEB
\end{TerminalOutput}
依据：\path|reports/deployment/container-checks.log|。

\par\noindent\textbf{终端输出（格式示意）}\par
\begin{TerminalOutput}
{"status":"PASS","records":"/mnt/robot_disk/remote_sensing_project/outputs/deployment-<本次目录>"}
\end{TerminalOutput}
示意中的任务名、时间、目录、哈希和耗时以实际运行值为准；该输出用于说明格式，不构成本次重新运行证据。

\par\noindent\textbf{完成判定与异常处理}\par
各命令均退出 0，相关检查输出 PASS；unittest 的跳过项需单独记录。完整部署测试中的 32 MiB OOM 是普通内存异常测试，不是 GPU 显存测试。

\section{TEST-02 Docker GPU 计算}
\par\noindent\textbf{用途}\par
验证容器内实际 CUDA 运算与 CPU 参考结果一致。

\par\noindent\textbf{执行前提}\par
宿主驱动、Toolkit 和 Docker GPU 运行时已配置；准备好本地镜像。

\par\noindent\textbf{需填写内容}\par
IMAGE：待测算法镜像；脚本缺省 0.3.0-kmeans1。本测试无输入影像和业务参数。

\par\noindent\textbf{执行命令}\par
\begin{Verbatim}
nvidia-smi
nvidia-ctk --version
bash scripts/test-gpu.sh
\end{Verbatim}
\par\noindent\textbf{终端输出（实测记录节选）}\par
\begin{TerminalOutput}
{
  "status": "PASS",
  "device": "NVIDIA GeForce GTX 1650",
  "operation": "float32 vector addition",
  "backend": "CUDA Driver API + PTX JIT",
  "cpu_fallback": false,
  "device_allocation_bytes": 12582996
}
\end{TerminalOutput}
依据：\path|reports/deployment/gpu-followup/compute-result.json|。仅列与判定有关的字段，省略内容不表示程序只输出这些字段。

\par\noindent\textbf{终端输出（实测记录节选）}\par
\begin{TerminalOutput}
"runs": [
  {"seed":0,"elements":1048583,"max_absolute_error":0.0},
  {"seed":1,"elements":1048583,"max_absolute_error":0.0},
  {"seed":2,"elements":1048583,"max_absolute_error":0.0}
]
\end{TerminalOutput}
依据：\path|reports/deployment/gpu-followup/compute-result.json|。仅列与判定有关的字段，省略内容不表示程序只输出这些字段。

\par\noindent\textbf{完成判定与异常处理}\par
status 为 PASS、\path|cpu_fallback| 为 false，三轮 \path|max_absolute_error| 均为 0。nvidia-smi 单独成功不满足该判定。该测试分配约 12 MiB 显存，不证明真实模型的 CUDA 框架与显存需求已验收。

\chapter{Kubernetes 操作与 GPU 调度}
\section{K8S-01 检查节点与导入算法镜像}
\par\noindent\textbf{用途}\par
确认目标集群并向 kind 节点加载已有 Docker 镜像。

\par\noindent\textbf{执行前提}\par
执行 ENV-01；本机 rs-lab 已存在。

\par\noindent\textbf{需填写内容}\par
IMAGE：已验收镜像；集群名称固定 rs-lab，目标集群另行适配。

\par\noindent\textbf{执行命令}\par
\begin{Verbatim}
kubectl config current-context
kubectl get nodes -o wide
export TMPDIR=/mnt/robot_disk/k8s/cache
kind load docker-image "$IMAGE" --name rs-lab
\end{Verbatim}
\begin{center}\small
\begin{tabular}{>{\raggedright\arraybackslash}p{5.62cm}>{\raggedright\arraybackslash}p{9.98cm}}
\toprule
\textbf{指令 / 参数} & \textbf{用途} \\ \midrule
current-context & 确认当前操作上下文，不切换集群 \\
-o wide & 显示节点地址、运行时等扩展字段 \\
kind load docker-image & 将 Docker 镜像加载到指定 kind 节点 \\
--name rs-lab & 指定本地集群 \\
TMPDIR & 导入临时文件使用数据盘缓存 \\
\bottomrule
\end{tabular}
\end{center}
\par\noindent\textbf{终端输出（格式示意）}\par
\begin{TerminalOutput}
kind-rs-lab
NAME                   STATUS   ROLES           VERSION
rs-lab-control-plane   Ready    control-plane   v1.36.4
Image ... loading...
\end{TerminalOutput}
示意中的任务名、时间、目录、哈希和耗时以实际运行值为准；该输出用于说明格式，不构成本次重新运行证据。

\par\noindent\textbf{完成判定与异常处理}\par
节点 Ready，镜像导入命令退出 0；镜像已存在时可能提示跳过。Docker 镜像库与节点 containerd 镜像库互相独立。

\section{K8S-02 CPU Job 回归}
\par\noindent\textbf{用途}\par
验证两个 CPU Job 并发运行及大样本 Job 的暂停取消。

\par\noindent\textbf{执行前提}\par
小样本已生成；IMAGE 已导入；本机 hostPath 映射已建立。

\par\noindent\textbf{需填写内容}\par
大样本目录固定 \path|data/input/kmeans-large|；边长 1000 对应 1000000 像元，不超过脚本参数缺省上限。

\par\noindent\textbf{执行命令}\par
\begin{Verbatim}
kubectl config set-context --current --namespace=default
mkdir -p data/input/kmeans-large
docker run --rm --network none \
  --user "$(id -u):$(id -g)" \
  -v "$PWD/data/input/kmeans-large:/fixture" \
  -v "$PWD/tests:/tests:ro" \
  --entrypoint python "$IMAGE" \
  /tests/create_kmeans_fixture.py /fixture 1000
python3 scripts/test-k8s.py
\end{Verbatim}
\par\noindent\textbf{终端输出（格式示意）}\par
\begin{TerminalOutput}
Context "kind-rs-lab" modified.
{"status":"PASS","records":".../reports/deployment/k8s-<时间戳>",
 "jobs":["rs-example-<时间戳>-0","rs-example-<时间戳>-1"],
 "cancelled":"rs-cancel-<时间戳>"}
\end{TerminalOutput}
示意中的任务名、时间、目录、哈希和耗时以实际运行值为准；该输出用于说明格式，不构成本次重新运行证据。

\par\noindent\textbf{完成判定与异常处理}\par
status=PASS；小任务生成成功标记，取消任务无未清理结果。set-context 会修改本地 kubeconfig 的当前命名空间。结果位于 \path|outputs/rs-example-*|，证据位于 \path|reports/deployment/k8s-*|。

\section{K8S-03 GPU 完整调度测试}
\par\noindent\textbf{用途}\par
验证 GPU 资源登记、实际计算、排队、释放、超额申请及超时。

\par\noindent\textbf{执行前提}\par
本集群总可分配 GPU 为 1，当前没有其他 GPU 任务；插件已就绪。

\par\noindent\textbf{需填写内容}\par
IMAGE 可由环境变量指定；KUBECONFIG / KUBECTL 可覆盖本机工具路径。不要并发启动多个本测试副本。

\par\noindent\textbf{执行命令}\par
\begin{Verbatim}
kubectl get node rs-lab-control-plane \
  -o jsonpath='{.status.allocatable.nvidia\.com/gpu}{"\n"}'
kubectl -n kube-system get pods -l app=nvidia-device-plugin
python3 scripts/test-k8s-gpu.py
\end{Verbatim}
\par\noindent\textbf{终端输出（实测记录）}\par
\begin{TerminalOutput}
{
  "status": "PASS",
  "allocatable_gpus": 1,
  "gpu_correctness": true,
  "exclusive_scheduling": true,
  "excess_request_unschedulable": true,
  "no_request_no_device": true,
  "job_deadline_exceeded_verified": true,
  "release_to_completion_seconds": 11.306236156000523,
  "completed_job": "queued-1788868434",
  "records": "/mnt/robot_disk/remote_sensing_project/reports/k8s-gpu/run-1788868434"
}
\end{TerminalOutput}
依据：\path|reports/k8s-gpu/run-1788868434/summary.json|。

\begin{center}\small
\begin{tabular}{>{\raggedright\arraybackslash}p{7.49cm}>{\raggedright\arraybackslash}p{8.11cm}}
\toprule
\textbf{结果字段} & \textbf{完成含义} \\ \midrule
\path|gpu_correctness| & 实际 CUDA 计算正确 \\
\path|exclusive_scheduling| & 保留配额时排队，释放后完成 \\
\path|excess_request_unschedulable| & 请求两个 GPU 时无法调度 \\
\path|no_request_no_device| & 未申请 GPU 的 Pod 看不到设备 \\
\path|job_deadline_exceeded_verified| & 实际触发 Job 超时 \\
\path|release_to_completion_seconds| & 包含调度、启动、计算和轮询；不是纯调度延迟 \\
\bottomrule
\end{tabular}
\end{center}
\par\noindent\textbf{完成判定与异常处理}\par
status=PASS 且五个布尔检查均为 true。allocatable=1 表示容量，不表示实时空闲。测试先计算再空闲保留配额，不能证明运行中的长 CUDA 内核可即时无损停止。

\section{K8S-04 手动提交 GPU Job}
\par\noindent\textbf{用途}\par
单独提交可复现的 CUDA 计算任务。

\par\noindent\textbf{执行前提}\par
K8S-03 基础条件满足；已存在 \path|configs/gpu/compute-job.json|。

\par\noindent\textbf{需填写内容}\par
默认无需修改模板。若换镜像，应修改模板的 \path|spec.template.spec.containers|[0].image；仅 export IMAGE 不会改变该 JSON。

\par\noindent\textbf{执行命令}\par
\begin{Verbatim}
kubectl create namespace rs-gpu-tests --dry-run=client -o yaml | \
  kubectl apply -f -
kubectl -n rs-gpu-tests create configmap rs-cuda-compute \
  --from-file=verify_cuda_compute.py=tests/verify_cuda_compute.py \
  --dry-run=client -o yaml | kubectl apply -f -
JOB=$(kubectl create -f configs/gpu/compute-job.json \
  -o jsonpath='{.metadata.name}')
printf 'JOB=%s\n' "$JOB"
kubectl -n rs-gpu-tests wait --for=condition=complete \
  "job/$JOB" --timeout=120s
kubectl -n rs-gpu-tests logs "job/$JOB"
\end{Verbatim}
\begin{center}\small
\begin{tabular}{>{\raggedright\arraybackslash}p{7.18cm}>{\raggedright\arraybackslash}p{8.42cm}}
\toprule
\textbf{参数 / 字段} & \textbf{含义} \\ \midrule
-n rs-gpu-tests & 目标命名空间 \\
--dry-run=client -o yaml & 生成对象描述，经 apply 创建或更新 \\
--from-file=键=文件 & 把计算脚本加入 ConfigMap \\
generateName & 自动生成不冲突的 Job 名称 \\
--for=condition=complete & 等待 Job 完成；不是创建后立即返回成功 \\
--timeout=120s & 本地 wait 命令最长等待时间 \\
\bottomrule
\end{tabular}
\end{center}
\par\noindent\textbf{终端输出（格式示意）}\par
\begin{TerminalOutput}
namespace/rs-gpu-tests configured
configmap/rs-cuda-compute created
JOB=rs-cuda-check-9d229
job.batch/rs-cuda-check-9d229 condition met
{"status":"PASS","device":"NVIDIA GeForce GTX 1650", ...}
\end{TerminalOutput}
示意中的任务名、时间、目录、哈希和耗时以实际运行值为准；该输出用于说明格式，不构成本次重新运行证据。

创建/配置提示依赖对象是否已存在。实际独立运算证据保存于 \path|reports/k8s-gpu/manual-job.log|。此处 Job 名以实际生成值为准，JSON 省略号仅用于展示。

\section{手动 GPU Job 模板参数}
\begin{center}\small
\begin{tabular}{>{\raggedright\arraybackslash}p{6.86cm}>{\raggedright\arraybackslash}p{3.43cm}>{\raggedright\arraybackslash}p{5.3cm}}
\toprule
\textbf{字段} & \textbf{当前值} & \textbf{修改规则} \\ \midrule
image / imagePullPolicy & 0.3.0-kmeans1 / Never & 修改镜像后先导入节点；Never 不在线拉取 \\
\path|requests.nvidia.com/gpu| & 1 & 本机只有 1；多卡目标另行设计 \\
\path|limits.nvidia.com/gpu| & 1 & 与申请数量一致 \\
requests.cpu / memory & 100m / 128Mi & 调度请求 \\
limits.cpu / memory & 1 / 512Mi & 容器 CPU / RAM 上限，不是显存 \\
activeDeadlineSeconds & 120 & Job 生命周期时限 \\
\path|terminationGracePeriodSeconds| & 10 & Pod 正常终止宽限 \\
backoffLimit & 0 & 失败不自动重试 \\
runAsUser / runAsGroup & 1000 / 1000 & 模板固定值；不是动态读取当前 UID \\
\bottomrule
\end{tabular}
\end{center}
普通计算 Pod 使用默认 runc 与 CDI 注入，不需要 privileged、显卡 hostPath 或 nvidia-bootstrap RuntimeClass。宿主桌面和外部 Docker 程序仍可竞争显存；Kubernetes 单卡配额不构成整卡性能或显存隔离。

\section{K8S-05 查看、暂停与保留 Job 记录}
\par\noindent\textbf{用途}\par
保存任务日志并暂停整个 Job，避免控制器重建被单独删除的 Pod。

\par\noindent\textbf{执行前提}\par
已知实际 Job 名称及其命名空间。

\par\noindent\textbf{需填写内容}\par
NS：命名空间；JOB：K8S-04 返回的名称，或查询所得 CPU Job 名称。

\par\noindent\textbf{执行命令}\par
\begin{Verbatim}
NS=rs-gpu-tests
# 本终端沿用 K8S-04 的 JOB；新终端须先填写 JOB=实际名称。
kubectl -n "$NS" get jobs,pods -o wide
kubectl -n "$NS" describe "job/$JOB"
mkdir -p logs
kubectl -n "$NS" logs "job/$JOB" > "logs/$JOB.log"
kubectl -n "$NS" patch "job/$JOB" --type=merge \
  -p '{"spec":{"suspend":true}}'
kubectl -n "$NS" get "job/$JOB" \
  -o jsonpath='{.spec.suspend}{"\n"}'
\end{Verbatim}
\par\noindent\textbf{终端输出（格式示意）}\par
\begin{TerminalOutput}
job.batch/rs-cuda-check-9d229 patched
true
\end{TerminalOutput}
示意中的任务名、时间、目录、哈希和耗时以实际运行值为准；该输出用于说明格式，不构成本次重新运行证据。

\par\noindent\textbf{完成判定与异常处理}\par
suspend=true 表示暂停设置生效，活动 Pod 的删除需要等待终止过程；已完成 Job 无运行进程可停止。恢复会重建 Pod，不等于算法断点续算。归档后只清理明确的测试对象。

\chapter{本地集群维护}
\section{MAINT-01 创建缺失的 kind 集群}
\par\noindent\textbf{用途}\par
在指定数据盘上建立 rs-lab。

\par\noindent\textbf{执行前提}\par
仅在 rs-lab 不存在时使用；kind、kubectl 已安装，Docker 数据目录位于指定数据盘。

\par\noindent\textbf{需填写内容}\par
当前脚本固定 /\path|mnt/robot_disk/k8s| 与 rs-lab；不提供通用安装参数。

\par\noindent\textbf{执行命令}\par
\begin{Verbatim}
bash scripts/create-kind.sh
\end{Verbatim}
\par\noindent\textbf{终端输出（格式示意）}\par
\begin{TerminalOutput}
Creating cluster "rs-lab" ...
Set kubectl context to "kind-rs-lab"
\end{TerminalOutput}
示意中的任务名、时间、目录、哈希和耗时以实际运行值为准；该输出用于说明格式，不构成本次重新运行证据。

\par\noindent\textbf{完成判定与异常处理}\par
以 kubectl get nodes 显示 Ready 为准。集群已存在时脚本会保留配置并退出 2，提示 rs-lab 已存在；不得将该情形解释为需要删除现有集群。

\section{MAINT-02 复装本机 GPU 接入}
\par\noindent\textbf{用途}\par
复制当前宿主驱动到 kind 节点，配置初始化运行时并安装设备插件。

\par\noindent\textbf{执行前提}\par
维护窗口内执行；没有活动或 Pending GPU Pod。宿主驱动和 Toolkit 已完成配置。

\par\noindent\textbf{需填写内容}\par
脚本固定本机节点和官方插件摘要；驱动升级、其他操作系统或生产集群应制定独立方案。

\par\noindent\textbf{执行命令}\par
\begin{Verbatim}
python3 scripts/setup-kind-gpu.py
bash scripts/install-kind-device-plugin.sh
python3 scripts/test-k8s-gpu.py
\end{Verbatim}
\par\noindent\textbf{终端输出（实测记录）}\par
\begin{TerminalOutput}
PASS: kind driver root, CDI and bootstrap runtime prepared
\end{TerminalOutput}
依据：\path|reports/k8s-gpu/setup-replay.log|。

\par\noindent\textbf{终端输出（实测记录节选）}\par
\begin{TerminalOutput}
daemon set "nvidia-device-plugin" successfully rolled out
PASS: nvidia.com/gpu=1
\end{TerminalOutput}
依据：\path|reports/k8s-gpu/install-replay.log|。仅列与判定有关的字段，省略内容不表示程序只输出这些字段。

\par\noindent\textbf{完成判定与异常处理}\par
准备、安装和完整复测均通过。准备脚本会删除自己的插件并重启节点内 containerd；不是日常运行前置步骤。首次缺少插件镜像需联网。驱动副本不是自动升级机制，不能直接照搬到生产节点。

\section{组件与路径}
\begin{center}\small
\begin{tabular}{>{\raggedright\arraybackslash}p{6.86cm}>{\raggedright\arraybackslash}p{8.74cm}}
\toprule
\textbf{对象} & \textbf{作用} \\ \midrule
/\path|opt/nvidia-driver| & 节点内驱动副本 \\
/\path|etc/cdi/nvidia-kind.json| & 插件初始化用 CDI 描述 \\
/\path|run/cdi/| & 设备插件生成的任务 CDI 描述 \\
nvidia-bootstrap & 只供设备插件初始化使用的 RuntimeClass \\
device-plugin DaemonSet & 向 kubelet 登记 \path|nvidia.com/gpu| \\
\bottomrule
\end{tabular}
\end{center}
install-docker.sh 和 install-docker-retry.sh 是本机首次安装历史脚本，包含用户和磁盘身份限制。当前 Docker 已存在，不应作为每次部署或目标机器的通用安装步骤执行。

\chapter{构建、安全扫描与离线交付}
\section{REL-01 构建新的镜像版本}
\par\noindent\textbf{用途}\par
固定基础镜像摘要和构建编号，生成独立算法镜像。

\par\noindent\textbf{执行前提}\par
构建网络可用，源码与依赖已审查；记录工作区修改。

\par\noindent\textbf{需填写内容}\par
PLATFORM：本手册为 \path|linux/amd64|；\path|BUILD_ID|：唯一编号；\path|PYTHON_BASE|：已核验摘要。

\par\noindent\textbf{执行命令}\par
\begin{Verbatim}
export PYTHON_BASE=python:3.13-slim-trixie@sha256:9d2e5553305c7c7b0097999bb17187c69b921ccd6bc9d40e4bb5ebe652c00285
export PLATFORM=linux/amd64
export BUILD_ID="$(git rev-parse --short HEAD)-$(date +%Y%m%d%H%M%S)"
bash scripts/build.sh
export IMAGE="remote-sensing-framework:$(cat VERSION)-$BUILD_ID"
printf 'IMAGE=%s\n' "$IMAGE"
\end{Verbatim}
\begin{center}\small
\begin{tabular}{>{\raggedright\arraybackslash}p{5.15cm}>{\raggedright\arraybackslash}p{10.45cm}}
\toprule
\textbf{变量} & \textbf{要求} \\ \midrule
\path|PYTHON_BASE| & 必须包含 @sha256: 后接 64 位十六进制摘要 \\
\path|BUILD_ID| & 字母或数字开头，仅允许字母、数字、下划线、点、连字符 \\
PLATFORM & 明确指定目标架构，不将本机架构视为甲方要求 \\
\path|IMAGE_REPOSITORY| & 可选；缺省 \path|remote-sensing-framework| \\
\bottomrule
\end{tabular}
\end{center}
\par\noindent\textbf{终端输出（格式示意）}\par
\begin{TerminalOutput}
... 构建过程输出 ...
remote-sensing-framework:0.3.0-e59ce62-20260908150000
IMAGE=remote-sensing-framework:0.3.0-e59ce62-20260908150000
\end{TerminalOutput}
示意中的任务名、时间、目录、哈希和耗时以实际运行值为准；该输出用于说明格式，不构成本次重新运行证据。

\par\noindent\textbf{完成判定与异常处理}\par
构建命令退出 0，并对新 IMAGE 执行 TEST-01。当前构建上下文包含工作区文件，提交号本身不能证明未提交源码未参与构建。该基础摘要是当前项目基线，不能视为永久安全版本。

\section{REL-02 扫描镜像漏洞}
\par\noindent\textbf{用途}\par
生成 \path|HIGH/CRITICAL| 漏洞 JSON，并用退出码实施门禁。

\par\noindent\textbf{执行前提}\par
Trivy 已安装；在线扫描需要数据库更新网络。

\par\noindent\textbf{需填写内容}\par
IMAGE：待扫描镜像；TRIVY：实际可执行文件；\path|SCAN_DIR|：本次独立报告目录。

\par\noindent\textbf{执行命令}\par
\begin{Verbatim}
export TRIVY=/mnt/robot_disk/tools/trivy-0.74.0/trivy
export SCAN_DIR="$PWD/reports/security/scan-$(date +%Y%m%d-%H%M%S)"
if bash scripts/scan-image.sh; then
  scan_rc=0
else
  scan_rc=$?
fi
printf 'scan_exit=%s\n' "$scan_rc"
ls "$SCAN_DIR"
\end{Verbatim}
\par\noindent\textbf{终端输出（格式示意）}\par
\begin{TerminalOutput}
scan_exit=2
scan.json
tool-version.txt
\end{TerminalOutput}
示意中的任务名、时间、目录、哈希和耗时以实际运行值为准；该输出用于说明格式，不构成本次重新运行证据。

\par\noindent\textbf{完成判定与异常处理}\par
0 表示本次所选级别未命中；2 表示高危门禁命中；其他非零需检查工具或网络故障。文件存在并不代表扫描通过。复用旧数据库时必须保留数据库时间。

\section{已记录安全状态}
\begin{center}\small
\begin{tabular}{>{\raggedright\arraybackslash}p{7.64cm}>{\raggedright\arraybackslash}p{3.74cm}>{\raggedright\arraybackslash}p{4.21cm}}
\toprule
\textbf{组件} & \textbf{HIGH / CRITICAL} & \textbf{结论} \\ \midrule
算法镜像 0.3.0-kmeans1 & 47 / 3 & 尚未达到安全门禁 \\
设备插件 v0.20.0 & 2 / 0 & 尚未达到安全门禁 \\
\bottomrule
\end{tabular}
\end{center}
算法镜像为 50 条软件包-漏洞记录、18 个不同编号。插件两条均为 CVE-2026-14456，扫描器列出了修复版本。以上来自同日历史扫描，本次文档修订未重新扫描或修复。

\section{REL-03 导出单个镜像}
\par\noindent\textbf{用途}\par
生成可 docker load 的压缩镜像及 SHA256 文件。

\par\noindent\textbf{执行前提}\par
IMAGE 已完成相应验收；导出目标有足够空间。

\par\noindent\textbf{需填写内容}\par
\path|EXPORT_FILE|：新的导出文件路径，现有文件会被拒绝覆盖。

\par\noindent\textbf{执行命令}\par
\begin{Verbatim}
export EXPORT_FILE="/mnt/robot_disk/deliveries/image-$(date +%Y%m%d-%H%M%S).tar.gz"
bash scripts/export.sh
sha256sum -c "$EXPORT_FILE.sha256"
\end{Verbatim}
\par\noindent\textbf{终端输出（格式示意）}\par
\begin{TerminalOutput}
/mnt/robot_disk/deliveries/image-20260908-150500.tar.gz: OK
\end{TerminalOutput}
示意中的任务名、时间、目录、哈希和耗时以实际运行值为准；该输出用于说明格式，不构成本次重新运行证据。

\par\noindent\textbf{完成判定与异常处理}\par
export.sh 通常无成功提示；以退出 0、两个文件存在和校验 OK 为准。该产物仅含镜像，不包含配置、输入数据、GPU 驱动或运维工具。

\section{REL-04 加载已有 CPU 离线交付包}
\par\noindent\textbf{用途}\par
验证外层及内部校验文件后离线加载镜像。

\par\noindent\textbf{执行前提}\par
目标为 Linux amd64，已有 Docker、Bash、Python 3、gzip 和 sha256sum。

\par\noindent\textbf{需填写内容}\par
\path|BUNDLE_DIR|：压缩包所在目录；\path|PACKAGE_ROOT|：解包后 SHA256SUMS 所在目录。

\par\noindent\textbf{执行命令}\par
\begin{Verbatim}
BUNDLE_DIR=/mnt/robot_disk/deliveries
cd "$BUNDLE_DIR"
sha256sum -c rs-0.3.0-kmeans1-bundle.tar.gz.sha256
EXTRACT_DIR="/mnt/robot_disk/tmp/rs-load-$(date +%Y%m%d-%H%M%S)"
mkdir -p "$EXTRACT_DIR"
tar -xzf rs-0.3.0-kmeans1-bundle.tar.gz -C "$EXTRACT_DIR"
find "$EXTRACT_DIR" -name SHA256SUMS -print
\end{Verbatim}
\par\noindent\textbf{终端输出（格式示意）}\par
\begin{TerminalOutput}
rs-0.3.0-kmeans1-bundle.tar.gz: OK
/mnt/robot_disk/tmp/rs-load-20260908-151000/rs-0.3.0-kmeans1/SHA256SUMS
\end{TerminalOutput}
示意中的任务名、时间、目录、哈希和耗时以实际运行值为准；该输出用于说明格式，不构成本次重新运行证据。

将 find 返回路径末尾的 /SHA256SUMS 去掉，填写为 \path|PACKAGE_ROOT|。以下示例目录须按实际结果替换，不假定压缩包内部目录固定。

\par\noindent\textbf{执行命令}\par
\begin{Verbatim}
PACKAGE_ROOT=/mnt/robot_disk/tmp/rs-load-20260908-151000/rs-0.3.0-kmeans1
cd "$PACKAGE_ROOT"
sha256sum -c SHA256SUMS
gzip -dc image.tar.gz | docker load
test -e .env || cp .env.example .env
\end{Verbatim}
\begin{center}\small
\begin{tabular}{>{\raggedright\arraybackslash}p{5.46cm}>{\raggedright\arraybackslash}p{10.14cm}}
\toprule
\textbf{参数} & \textbf{含义} \\ \midrule
sha256sum -c & 按清单校验文件；必须全部 OK \\
tar -xzf ... -C & 解压 gzip tar 到指定新目录 \\
gzip -dc & 解压到 stdout，不删除原压缩包 \\
docker load & 从标准输入加载镜像 \\
test -e .env || cp & 仅在 .env 不存在时复制，避免覆盖用户配置 \\
\bottomrule
\end{tabular}
\end{center}
\par\noindent\textbf{终端输出（格式示意）}\par
\begin{TerminalOutput}
image.tar.gz: OK
... 其他交付文件: OK
Loaded image: remote-sensing-framework:0.3.0-kmeans1
\end{TerminalOutput}
示意中的任务名、时间、目录、哈希和耗时以实际运行值为准；该输出用于说明格式，不构成本次重新运行证据。

\par\noindent\textbf{完成判定与异常处理}\par
全部内部校验通过且镜像成功加载后，按配置及执行章节运行。若在解包目录工作，\path|PROJECT_ROOT| 应设为该目录。离线 CPU 包不包含完整 \path|GPU/K8s| 基础设施；旧包也不自动包含本修订手册。

\section{REL-05 使用现有打包入口}
\par\noindent\textbf{用途}\par
生成当前脚本定义的交付目录，包含镜像、样本、源码及指定文档。

\par\noindent\textbf{执行前提}\par
IMAGE 已验证；输出目录尚不存在。正式完整交付前需补齐文档清单。

\par\noindent\textbf{需填写内容}\par
\path|PACKAGE_DIR|：新目录绝对路径。

\par\noindent\textbf{执行命令}\par
\begin{Verbatim}
cd "$PROJECT_ROOT"
export PACKAGE_DIR="/mnt/robot_disk/deliveries/rs-package-$(date +%Y%m%d-%H%M%S)"
bash scripts/package.sh
\end{Verbatim}
\par\noindent\textbf{终端输出（格式示意）}\par
\begin{TerminalOutput}
交付目录：/mnt/robot_disk/deliveries/rs-package-20260908-152000
\end{TerminalOutput}
示意中的任务名、时间、目录、哈希和耗时以实际运行值为准；该输出用于说明格式，不构成本次重新运行证据。

\par\noindent\textbf{完成判定与异常处理}\par
此命令生成目录，不自动生成外层 bundle.tar.gz。当前 package.sh 固定文档清单未包含本修订手册和新 GPU 总报告；不得将输出直接标为文档齐全的新交付版。增补文件后必须重建并校验 SHA256SUMS。本次未修改打包脚本或旧包。

\chapter{源码推送与文件发布}
\section{GIT-01 推送已提交开发分支}
\par\noindent\textbf{用途}\par
将已有本地提交上传到 origin 的开发分支。

\par\noindent\textbf{执行前提}\par
仓库 SSH 凭据可用；所需文件已经明确提交。

\par\noindent\textbf{需填写内容}\par
分支为 \path|codex/kmeans-deployment|；如分支变更，应先核对再调整。

\par\noindent\textbf{执行命令}\par
\begin{Verbatim}
cd /mnt/robot_disk/remote_sensing_project
git status --short
git branch --show-current
git log -3 --oneline
git push -u origin codex/kmeans-deployment
\end{Verbatim}
\begin{center}\small
\begin{tabular}{>{\raggedright\arraybackslash}p{6.08cm}>{\raggedright\arraybackslash}p{9.52cm}}
\toprule
\textbf{参数} & \textbf{含义} \\ \midrule
status --short & 检查尚未提交的修改 \\
branch --show-current & 确认当前分支 \\
log -3 --oneline & 确认最近三个提交 \\
push -u origin 分支 & 推送该分支并设置上游跟踪；不合并 main \\
\bottomrule
\end{tabular}
\end{center}
\par\noindent\textbf{终端输出（格式示意）}\par
\begin{TerminalOutput}
To github.com:huam54925-cpu/remote-sensing-algorithms.git
 * [new branch]      codex/kmeans-deployment -> codex/kmeans-deployment
branch ... set up to track ...
\end{TerminalOutput}
示意中的任务名、时间、目录、哈希和耗时以实际运行值为准；该输出用于说明格式，不构成本次重新运行证据。

\par\noindent\textbf{完成判定与异常处理}\par
实际输出也可能是增量更新或 Everything up-to-date。git push 不上传未提交手册，也不上传 Docker 镜像。新手册需要单独审阅、提交后才可推送；本次仅生成文档，不自动提交或发布。

\chapter{故障定位与交付验收}
\section{故障对照表}
\begin{center}\small
\begin{tabular}{>{\raggedright\arraybackslash}p{4.84cm}>{\raggedright\arraybackslash}p{5.3cm}>{\raggedright\arraybackslash}p{5.46cm}}
\toprule
\textbf{可见现象} & \textbf{检查项目} & \textbf{操作结论} \\ \midrule
输入或输出文件错误 & 目录权限、\path|RS_INPUT|、只读挂载 & 纠正路径/权限后采用新任务目录 \\
result 已存在 & \path|OUTPUT_DIR/result| & 换新输出目录，不覆盖旧成果 \\
退出 137 & OOMKilled、终止记录、日志 & 区分 OOM 与强杀 \\
GPU 不可见 & 宿主 nvidia-smi、Docker GPU 测试 & 逐层定位，不直接修改算法 \\
Insufficient \path|nvidia.com/gpu| & 申请量与当前配额占用 & 等待或减少请求；本机总量为 1 \\
ErrImageNeverPull & 节点镜像与模板 image & 导入完全一致的镜像 \\
CDI mount 错误 & Pod 事件、节点驱动文件 & 维护窗口检查 CDI 引用 \\
DeadlineExceeded & Job 时限、任务进展 & 确认是正常超时、资源不足或程序卡住 \\
\bottomrule
\end{tabular}
\end{center}
\section{DIAG-01 保存集群故障记录}
\par\noindent\textbf{用途}\par
收集排错所需状态，保留可复现信息。

\par\noindent\textbf{执行前提}\par
已设置 KUBECONFIG；目标命名空间存在。

\par\noindent\textbf{需填写内容}\par
NS：实际命名空间；RECORD 自动生成。单任务应另保存其日志、env 和退出码，外发前去除凭据。

\par\noindent\textbf{执行命令}\par
\begin{Verbatim}
NS=rs-gpu-tests
RECORD="logs/diagnosis-$(date +%Y%m%d-%H%M%S)-$$"
mkdir -p "$RECORD"
date -Is > "$RECORD/time.txt"
docker version > "$RECORD/docker-version.txt" 2>&1
findmnt -T "$PROJECT_ROOT" > "$RECORD/mount.txt"
nvidia-smi > "$RECORD/gpu.txt" 2>&1
kubectl get nodes -o wide > "$RECORD/nodes.txt" 2>&1
kubectl -n "$NS" get jobs,pods -o wide \
  > "$RECORD/tasks.txt" 2>&1
kubectl -n "$NS" get events \
  --sort-by=.metadata.creationTimestamp \
  > "$RECORD/events.txt" 2>&1
printf 'RECORD=%s\n' "$RECORD"
\end{Verbatim}
\par\noindent\textbf{终端输出（格式示意）}\par
\begin{TerminalOutput}
RECORD=logs/diagnosis-20260908-153000-12345
\end{TerminalOutput}
示意中的任务名、时间、目录、哈希和耗时以实际运行值为准；该输出用于说明格式，不构成本次重新运行证据。

\par\noindent\textbf{完成判定与异常处理}\par
目录创建成功不等于所有诊断命令成功；逐个查看文件，命令错误也可能被重定向保存。不要公开 kubeconfig、私钥或完整环境变量。容器 --rm 后的日志应依靠启动时保存的宿主日志。

\section{交付完成条件}
\begin{itemize}
\item 配置：必填字段明确，样例能运行；宿主与容器路径对应。
\item 结果：退出码、成功标记和期望文件一致；真实数据按对方基准比较。
\item 生命周期：启动、日志、停止、异常和重跑行为均有验证记录。
\item 资源：\path|CPU/RAM/GPU| 需求经过目标负载验证；GPU 诊断不替代真实框架测试。
\item 部署：干净机器或目标服务器完成独立离线演练。
\item 安全：高危门禁满足要求，组件、许可证、权重和 Secrets 处理明确。
\item 交付物：镜像、env、命令手册、数据及校验清单版本一致。
\end{itemize}
\section{待真实代码到达后填写的信息}
\begin{center}\small
\begin{tabular}{>{\raggedright\arraybackslash}p{4.21cm}>{\raggedright\arraybackslash}p{11.39cm}}
\toprule
\textbf{类别} & \textbf{应提供或确认的内容} \\ \midrule
入口与依赖 & 源程序入口、依赖锁定、框架/CUDA 版本、构建步骤 \\
输入 & 标准数据、波段、缩放、NoData、CRS、尺寸限制 \\
输出 & 标准结果、文件格式、误差容限、成功判定 \\
资源与停止 & 峰值内存/显存、运行时长、子进程、取消与检查点 \\
联网与服务 & 权重下载、授权、外部接口、是否需常驻 API \\
目标平台 & 架构、运行时、K8s、存储、镜像仓库和权限 \\
\bottomrule
\end{tabular}
\end{center}
当前尚未完成的重点包括漏洞修复、异机验收、磁盘满、大规模输入、并发写入、长 GPU 内核取消、子进程清理、日志留存策略、Kubernetes 网络策略与生产存储。该清单不将尚未实现的能力作为交付承诺。

\chapter{附录：脚本、证据与文档维护}
\section{脚本与实际作用}
\begin{center}\small
\begin{tabular}{>{\raggedright\arraybackslash}p{7.02cm}>{\raggedright\arraybackslash}p{8.58cm}}
\toprule
\textbf{脚本} & \textbf{实际作用} \\ \midrule
env-run.py → run.sh & 解析配置，组装挂载和资源限制，exec Docker \\
stop.sh & 调用 docker stop，传入任务名与宽限秒数 \\
test-gpu.sh & 调用已有 \path|verify_cuda_compute.py|，执行诊断内核 \\
test-k8s-gpu.py & 创建测试资源、检查调度与计算、保存证据 \\
setup-kind-gpu.py & 准备 kind 驱动副本、CDI 与专用运行时 \\
\path|install-kind-device-plugin.sh| & 导入固定插件镜像、部署并检查资源登记 \\
export.sh / package.sh & 镜像导出 / 交付目录制作 \\
inspect-algorithm.py & 只读搜索依赖及调用线索，不执行收到的代码 \\
\bottomrule
\end{tabular}
\end{center}
\section{日志和结果字段}
\begin{center}\small
\begin{tabular}{>{\raggedright\arraybackslash}p{6.08cm}>{\raggedright\arraybackslash}p{9.52cm}}
\toprule
\textbf{字段} & \textbf{解释} \\ \midrule
timestamp / level / module & 日志时间、级别、模块 \\
\path|task_id| / algorithm & 关联任务和算法 \\
message & \path|algorithm_started|、\path|algorithm_completed|、\path|task_cancelled| 等事件 \\
\path|elapsed_seconds| & 进程内部记录耗时，不等于完整平台延迟 \\
\path|process_peak_rss_kib| & 进程 RSS 峰值，不含完整 cgroup 或 GPU 显存 \\
\path|valid_pixels| / \path|invalid_pixels| & 有效与无效像元数 \\
status=completed & success.json 的执行完成状态 \\
\bottomrule
\end{tabular}
\end{center}
\section{实测证据位置}
\begin{itemize}
\item 总体状态：\path|docs/readiness-report-2026-09-08.md|。
\item CPU 演练及退出码：\path|reports/deployment/docker-integration/|。
\item Docker GPU：\path|reports/deployment/gpu-followup/compute-result.json|。
\item Kubernetes GPU：\path|reports/k8s-gpu/run-1788868434/summary.json|。
\item 手动 GPU Job：\path|reports/k8s-gpu/manual-job.log|。
\item 插件扫描：\path|reports/k8s-gpu/plugin-security.json|。
\end{itemize}
\section{文档源文件与编译}
本版提供 Markdown 文本版和 LaTeX 源码。Markdown 用于复制命令；LaTeX 用于排版 PDF。修订时应同步维护内容，不将终端输出示例复制为执行指令。

\par\noindent\textbf{执行命令}\par
\begin{Verbatim}
cd /mnt/robot_disk/remote_sensing_project
bash docs/manual/build.sh
sha256sum -c output/pdf/remote-sensing-operations-manual.pdf.sha256
\end{Verbatim}
\par\noindent\textbf{终端输出（格式示意）}\par
\begin{TerminalOutput}
PDF: /mnt/robot_disk/remote_sensing_project/output/pdf/remote-sensing-operations-manual.pdf
output/pdf/remote-sensing-operations-manual.pdf: OK
\end{TerminalOutput}
示意中的任务名、时间、目录、哈希和耗时以实际运行值为准；该输出用于说明格式，不构成本次重新运行证据。

编译使用 XeLaTeX，不执行文档内命令。本版按要求不逐页视觉检查，采用编译、字体缺字、文本、目录、命令块语法检查。代码块形式可复用不表示本次已执行其中每项维护操作。


\end{document}
